\section{Related Work}

\paragraph{Articulated-object manipulation}
Articulated object manipulation presents significant challenges due to diverse object geometries and physical characteristics. Although benchmarks by \citep{urakami2019doorgym, mu2021maniskill} have been introduced, their scope remains limited. While methods like \citep{peterson2000high, jain2009pulling, burget2013whole} have delved into motion planning, visual affordance learning has also gained attention \citep{mo2021where2act, wu2022vatmart, wang2022adaafford, zhao2022dualafford, geng2022end}. 
Other works \citep{eisner2022flowbot3d, xu2022universal} have proposed special representations for manipulation, though they are tailored mainly for suction grippers. 
On the other hand, \citep{geng2022gapartnet, geng2023partmanip} introduce generalizable parts as a novel representation but are limited to single-part interaction and can not handle natural instructions.
% Unique representations for this manipulation have been proposed by \citep{eisner2022flowbot3d, xu2022universal}, though they are tailored mainly for suction grippers. \citep{geng2022gapartnet, geng2023partmanip} introduce generalizable part as representation.



% Manipulating articulated objects is an important and open research area due to the complexity of different objects' geometry and physical properties. Previous work has proposed some benchmark \citep{urakami2019doorgym, mu2021maniskill} but the diversity is limited. As for the methods, \citep{peterson2000high, jain2009pulling, burget2013whole} explored motion planning, and \citep{mo2021where2act,wu2022vatmart,wang2022adaafford,zhao2022dualafford, geng2022end} leverages visual affordance learning. Other works \citep{eisner2022flowbot3d, xu2022universal} design special representations for articulated object manipulation and can generalize to a novel object category, but their methods are only suitable for the suction gripper. 


\paragraph{Generalizable object manipulation} Achieving generalization in robot applications is essential yet challenging. Various works \citep{fang2020graspnet, sundermeyer2021contact, gou2021rgb, gao2021kpam, xu2021adagrasp} employ supervised learning with motion planning to tackle generalizable tasks such as grasping. Nevertheless, these techniques often necessitate task-specific architectures and may falter in intricate manipulations. Though reinforcement learning holds promise for complex challenges \citep{rajeswaran2017dapg, akkaya2019solving}, the lack of generalizability remains unresolved \citep{kirk2021generalization_survey,ghosh2021generalization}. The ManiSkill benchmark \citep{mu2021maniskill} pioneers category-level object manipulation in simulations, while \citep{shen2022learning} utilizes imitation learning with RL-trained demonstrations. However, in both cases, the generalization ability is still unclear.
\citep{geng2022gapartnet, geng2023partmanip} first tackles generalizable manipulation in a cross-category manner, but is still limited to predefined part classes. \citep{xu2023unidexgrasp, wan2023unidexgrasp++} have universal generalization ability but only for grasping tasks. In summary, Generalizable object manipulation is still extremely challenging.

% For object manipulation, \citep{mu2021maniskill} challenges the acquisition of skills for articulated objects within familiar categories. The grander vision is perceiving and manipulating novel object categories, essential for numerous daily tasks. Despite robust algorithms for basic rigid objects \citep{fang2020graspnet,jiang2021synergies,breyer2020volumetric,li2018push,yu2016more}, limited research targets articulated objects with movable elements. Notably, \citep{mo2021where2act, wu2022vatmart} introduce low-level generalizability through per-pixel action estimation, exhibiting proficiency with unfamiliar objects. A closely related study by \citep{gadre2021act} introduces an interactive perception framework, though it overlooks consistent patterns for similar class parts and is restricted to simpler objects.

% Generalization is crucial yet challenging for robot application.
% Many works\citep{fang2020graspnet, sundermeyer2021contact, gou2021rgb, gao2021kpam, xu2021adagrasp} combine supervise learning and motion planning to learn generalizable skills, \textit{e.g.} grasping. However, these methods often require special architecture design for each task and may be unsuitable for complex manipulation tasks. Reinforcement learning (RL) has the potential to solve more complex task\citep{rajeswaran2017dapg, akkaya2019solving}, but the generalization ability of RL is an unsolved problem\citep{kirk2021generalization_survey,ghosh2021generalization}. 
% To facilitate the research of generalizable manipulation skills, ManiSkill \citep{mu2021maniskill} proposes the first category-level object manipulation benchmark in simulation.  \citep{shen2022learning} leverages imitation learning to learn complex generalizable manipulation skills from the demonstration. However, their demonstrations are collected by RL trained on every single instance, which requires a lot of effort to tune. In contrast, our experts can be directly trained on each object category. 


% On the front of object manipulation, Mu \textit{et al.} proposes \citep{mu2021maniskill} a challenge to learning generalizable manipulation skills for articulated objects from known categories.
% %
% Beyond category-level manner, one important further power is to generalize the perception and manipulation of complex objects in novel object categories, which is ubiquitously requested in handling many tasks in our daily lives.
% Although, for simple rigid objects, we have already seen robust and object-agnostic object grasping \citep{fang2020graspnet,jiang2021synergies,breyer2020volumetric} and planar pushing \citep{li2018push,yu2016more} algorithms, very few works have been devoted to interacting with articulated objects that contain movable parts. 
% Recently, Mo \textit{et al.}\citep{mo2021where2act} and Wu \textit{et al.}   \citep{wu2022vatmart} tackled this problem leveraging a low-level generalizability -- they estimate per-pixel actionability (\textit{e.g.} pushing, pulling) or motion trajectories on the object surface and shows that their trained network can generalize to unknown objects. 
% The most related work to us is Gadre \textit{et al.} \citep{gadre2021act} which proposes an interesting interactive perception pipeline that first learns to touch the object, watch its movement, and then segment the object into movable parts. However, this work doesn't consider the consistent geometry and actionability patterns behind parts from the same class and can only deal with simple objects with up to three parts on the table surfaces, \textit{e.g. scissors, eyeglasses}.

\paragraph{Robot control with large language models}
%
Incorporating pre-trained language models into robot systems has been gaining tremendous attention recently. A notable segment of these works concentrates on the grounding of language to manipulation skills, essential for executing long-horizon instructions~\citep{jang2022bc,stepputtis2020language,mees2022matters,codevilla2018driving,shao2021concept2robot,akakzia2021grounding,gong2023arnold,goyal2021zero,shridhar2022cliport,hu2019hierarchical, you2023make}. A hierarchical approach, seen in another subset of studies, hinges on a two-fold process: establishing a skill set and subsequently strategizing over it using large language models (LLMs) \citep{brohan2023can,huang2022inner,singh2022progprompt,huang2022language,lin2023text2motion,jiang2022vima}. However, this method is limited by the huge data collection effort, is still constrained to the predefined skill sets, and lacks visual understanding. In contrast, LLMs are employed to spawn code for low-level skills through APIs, though predicting these skills' continuous parameters remains challenging \citep{codeaspolicies2022}.
Some employ a way-point-based method for efficiency, sidestepping low-level robot actions in favor of trajectory hand poses~\citep{shridhar2023perceiver, gong2023arnold}. But they also require lots of data and have no detailed planning strategy. VoxPoser
\citep{huang2023voxposer} performs motion planning based on affordance and obstacle maps generated by LLM, but still lacks visual understanding.
%
For articulated objects with diverse geometries, intricate structures, and physical constraints, language-guided manipulation remains a challenging and underexplored problem.

% In another perspective, LMs for embodied applications predominantly revolve around planning and reasoning, spanning textual scene descriptions and perception APIs to vision usage during decoding or as direct input for multi-modal LLMs~\citep{xie2023translating,lu2023multimodal,patel2023pretrained,jansen2020visually,wang2023voyager,huang2022inner,zeng2022socratic,singh2022progprompt,codeaspolicies2022,huang2023grounded,driess2023palm,openai2023gpt}. A crucial component for these embodied LMs is their ability to act, which is often underpinned by a catalog of preset primitives. While some works highlight LLMs' behavioral sense for low-level control, challenges persist, including the need for custom motion primitives and questions surrounding compositional capacities at the spatial level~\cite {codeaspolicies2022}.


